\textit{The premises.}\enspace The identification of a tower's gauge algebra with the automorphism algebra of its minimal massive hierarchy is H2 (\claimref{dark-tower-matter.h2-identification}); the definition of a constituent, its field and the two factors' groups is Premise U (\claimref{dark-tower-matter.u-transposition-rule}), clauses U1 to U3; in dimensions three and five the constituent's momentum is H1 (\claimref{dark-tower-matter.h1-momentum}), which in dimension four is T12's own setting. Nothing else is assumed. The six causal towers and their signatures are those of Theorem \claimref{dark-tower-matter.constituents-field}.

\emph{The inputs, tower by tower.}  The field $F$ is read from Theorem \claimref{dark-tower-matter.constituents-field}: $\C$ at $r = 6, 2$, $\R$ at $r = 1, 7$, $\HH$ at $r = 3, 5$. The check recomputes it as $f\Cl^{0}(V)f$ with $f = \tfrac12(1 + e_{i}e_{j})$, $Q(e_{i}) = 1$, $Q(e_{j}) = -1$, the field acting on the right of $S = \Cl^{0}(V)f$ as in Theorem \claimref{dark-tower-matter.ladder}: it has real dimension $1$, $2$ (commutative) and $4$ (not commutative) in the three pairs of towers.

The onset count $a$ is Theorem \claimref{dark-tower-matter.where-mass-begins}: in dimensions three and four the reversion pairing on $S$ is alternating, a single constituent is massless and mass begins at two; in dimension five it is not alternating and mass begins at one. The check recomputes the criterion: the symmetrised pairing $\rev{\psi}_{1}\psi_{2} + \rev{\psi}_{2}\psi_{1}$ vanishes on every pair of basis elements of $S$ exactly when $n \le 4$.

The pair count $b$ and the pairs' scalars are Proposition \claimref{dark-tower-matter.pairs}: the real Witt index of $\hatV$ is $2$ in every causal tower and the complex one is $\lfloor (n+2)/2 \rfloor$, so $b = 2$ with real pairs for $n = 3$ and $b = 3$, the third pair needing a complex unit, for $n = 4, 5$. By U3 that unit lies in $F$. For $F = \C$ (this tower and the mirror) it is fixed up to sign and its centraliser in $F$ is $F$ itself, so $K = \C$. For $F = \HH$ Theorem \claimref{dark-tower-matter.third-pairs-unit-five-dimensions}(i) and (iii) give the centraliser of the pair algebra in $\HH$ as $\mathrm{span}\{1, u\}$, so again $K = \C$; the check confirms in $f\Cl^{0}(V)f$ that the centraliser of a unit $u$ is two-dimensional, and that in the commutator space $[F, F] = \mathfrak{sp}(1)$, the slot algebra of one quaternionic constituent, it is one-dimensional, the circle of Theorem \claimref{dark-tower-matter.gauge-algebras}(ii). For $F = \R$ no unit is needed, and none would be available: the case in which the pairs need a unit and $F$ supplies none does not occur among the causal towers.

The slot group is the momentum-preserving group of $a$ constituents, $O(a)$, $U(a)$ or $Sp(a)$ as $F = \R, \C, \HH$ (Theorem \claimref{dark-tower-matter.ladder}), cut in the quaternionic case to the subgroup commuting with $u$, which for $a = 1$ is $U(1)$; the pair group is, by U2, the unitary group over $K$ of the $b$ pairs, $O(b)$ or $U(b)$.

\emph{The automorphism algebras.}  An automorphism of $K^{a} \otimes_{K} K^{b}$ is $A \otimes B$; $A \otimes B = 1$ forces $A = \lambda 1$, $B = \lambda^{-1}1$ with $\lambda$ a unit scalar of $K$ (the argument of Theorem \claimref{standard-model.gauge-algebra-minimal-massive-composite}, a linear map fixing every line being a scalar). Over $\C$ the kernel is a circle and $\mathfrak{g}_{r} = (\uu(a) \oplus \uu(b))/\uu(1) = \su(a) \oplus \su(b) \oplus \uu(1)$, of dimension $a^{2} + b^{2} - 1$, as in Proposition \claimref{standard-model.four-dimensions-forces-converse}(2); over $\R$ it is $\{\pm 1\}$, discrete, and $\mathfrak{g}_{r} = \so(a) \oplus \so(b)$, of dimension $\tfrac12 a(a-1) + \tfrac12 b(b-1)$. The check computes each $\mathfrak{g}_{r}$ as the image of $\mathfrak{g}(a, K) \oplus \mathfrak{g}(b, K)$ acting by $X \otimes 1 + 1 \otimes Y$ on the real form of $K^{a} \otimes K^{b}$, as an exact rank, for all $a, b \le 3$ over both fields; it confirms that every image element is antisymmetric (it preserves the norm) and, over $\C$, commutes with the complex structure, and it computes the centre and the derived algebra: for $\C^{2} \otimes \C^{3}$ dimension $12$, centre $1$, derived algebra $11 = 3 + 8$, the images of $\su(2)$ and $\su(3)$; for $\C^{3} \otimes \C^{3}$, $17$, $1$, $16$; for $\C^{1} \otimes \C^{3}$, $9$, $1$, $8$; for $\R^{2} \otimes \R^{2}$, $2$, abelian. Substituting the inputs gives the table.

\emph{(i) The four-dimensional row.}  At $r = 6$ the inputs are T12's. Two constituents: Theorem \claimref{dark-tower-matter.where-mass-begins}(i) in dimension four is T11's two-constituent formula (Proposition \claimref{mass.mass-symplectic-area-two-twistor-formula}) written in the even algebra, the pairing $\rev{\psi}_{1}\psi_{2}$ in a copy of $\C$ playing the bracket (Remark \claimref{note.dark-tower-matter.four-dimensional-case-t11s-formula}). Three pairs: the complex Witt index $3$ of Proposition \claimref{dark-tower-matter.pairs}(ii) is Lemma \claimref{standard-model.pair-count-dimension} at $n = 4$. The field and the groups: $F = \C$, and by T14c's Remark \claimref{note.dark-tower-matter.why-u} any complex unit commuting with the spin action, the complexification's included, acts on the constituent ideal as $\pm I_{6}$, so T14c's real construction and T12's complexified one give the same $\C$; the slot group $U(2)$ of Theorem \claimref{dark-tower-matter.ladder} is T12's $U(2)$ (Proposition \claimref{standard-model.internal-symmetry-ladder}), and U2 takes the pair group $U(3)$ as T12 does. So the rule's hierarchy is $\C^{2} \otimes \C^{3}$, with automorphism algebra $\su(2) \oplus \su(3) \oplus \uu(1)$ of dimension $12$: Theorem \claimref{standard-model.gauge-algebra-minimal-massive-composite}. At $a = 3$ the same computation gives its seventeen-dimensional control. The rows $r = 1, 7$ (dimension $2$) and $r = 3, 5$ (dimension $9$) are not the Standard Model algebra, by Proposition \claimref{standard-model.four-dimensions-forces-converse}(2) for the complex rows and by dimension for the real ones.

\emph{(ii) The mirror row.}  At $r = 2$, signature $(3,1)$ and conformal signature $(4,2)$, every input is the same as at $r = 6$. The field is $\C$ by Theorem \claimref{dark-tower-matter.constituents-field}; the check confirms that the centre of $\Cl^{0}(4,2)$, like that of $\Cl^{0}(2,4)$, is $\mathrm{span}\{1, I_{6}\}$ with $I_{6}^{2} = -1$, and that $e_{K} \mapsto (-1)^{|K|/2}e_{K}$ is an algebra isomorphism of the even algebras in both the tower and the conformal signatures (Theorem \claimref{dark-tower-matter.mirror-algebra}(i)). The onset count is two (dimension four), the pairs are three and complex (Proposition \claimref{dark-tower-matter.pairs}(ii) names the mirror), and the ladder is $U(n)$ (Theorem \claimref{dark-tower-matter.ladder} names the mirror). So the mirror's row is this tower's, $\su(2) \oplus \su(3) \oplus \uu(1)$, which is the gauge algebra of the mirror copy theorem, Theorem \claimref{dark-matter.mirror-tower-copy-one}(i). The two derivations are independent: T14 runs T12's chain on the reversed data over $V_{\C}$, T14c's rule reads the real even algebras. The agreement uses U1. The commutant of the full Clifford module, computed in the check with a primitive idempotent, is $\HH$ for $\Cl(1,3)$ and $\R$ for $\Cl(3,1)$: with that reading the mirror's pairs would need a complex unit that $\R$ does not supply, and the two rows would not agree, which is T14c's reason (Remark \claimref{note.dark-tower-matter.why-u}) for rejecting it; that commutant is T14's ground type, a different object from the constituent field (Remark \claimref{note.dark-tower-matter.t14s-types-are-different-commutant}). Theorem \claimref{dark-tower-matter.mirror-algebra}(ii) sends the mirror's conformal unit to $-I_{6}$: the mirror's $\C$ is ours with the opposite orientation, which conjugates the representations and leaves the algebra unchanged.

\emph{(iii) The odd rows}  are Theorem \claimref{dark-tower-matter.gauge-algebras}: (i) for $r = 1, 7$ and (ii) for $r = 3, 5$.

\emph{(iv) $F$ and $K$.}  In the five-dimensional towers $F = \HH$ and $K = \C$, by the centraliser computation above. A rule that took $K = F$ there would keep the quaternionic slot group: that is the external-unit branch of T14c's Remark \claimref{note.dark-tower-matter.branches-break-u-what-they-give}, which breaks U3 and is not this rule.

\emph{What fails with wrong inputs.}  The table is sensitive to each input: three constituents in dimension four give $17$, two pairs give $7$, real pairs ($\R^{2} \otimes \R^{3}$) give $4$, one constituent gives $9$; none is $12$. In the check, replacing the onset count, the pair count, the field column or the scalars rule by another value makes it fail.
